% Original analytic teaching diagram, 2026-11-12.
% Quantization pipeline: source value x -> encoder alpha (cell index j) -> optional lossless index code
% (bit string) -> index decoder -> reconstruction c_j. Example values from the two-point Laplace quantizer:
% x = 0.7 -> j = 2 -> c_2 = 1.
\definecolor{VB}{HTML}{7188A3}\definecolor{VBF}{HTML}{E5EAF0}
\definecolor{VG}{HTML}{829C8B}\definecolor{VGF}{HTML}{E8EEE9}
\definecolor{VY}{HTML}{B4A16F}\definecolor{VYF}{HTML}{F1ECDD}
\definecolor{VR}{HTML}{A7787D}\definecolor{VRF}{HTML}{F0E5E6}
\definecolor{VP}{HTML}{9588A5}\definecolor{VPF}{HTML}{ECE8F0}
\definecolor{VI}{HTML}{3B4252}\definecolor{VM}{HTML}{6B7280}\definecolor{VL}{HTML}{4C4D4F}
\begin{tikzpicture}[x=.92mm,y=.92mm,>=Stealth,font=\figurefont\fontsize{8}{10}\selectfont,
lab/.style={text=VM,font=\figurefont\fontsize{7.5}{9}\selectfont},
box/.style={draw=VL,line width=.6pt,minimum width=16mm,minimum height=9mm,align=center,text=VI},
opt/.style={draw=VL,dashed,line width=.6pt,minimum width=16mm,minimum height=9mm,align=center,text=VI},
ar/.style={->,draw=VL,line width=.6pt}]
\node[text=VI] (x) at (0,0) {$x$};
\node[box,fill=VBF] (enc) at (16,0) {编码器\\$\alpha$};
\node[opt] (lc) at (41,0) {索引编码};
\node[opt] (ld) at (71,0) {索引解码};
\node[box,fill=VGF] (dec) at (96,0) {解码器\\$\beta$};
\node[text=VI] (y) at (112,0) {$\hat{x}$};
\draw[ar] (x) -- (enc);
\draw[ar] (enc) -- node[above,text=VI] {$j$} (lc);
\draw[ar] (lc) -- node[above,lab] {比特串} (ld);
\draw[ar] (ld) -- node[above,text=VI] {$j$} (dec);
\draw[ar] (dec) -- (y);
\node[lab] at (0,-9) {$0.7$};
\node[lab] at (28.5,-9) {$2$};
\node[lab] at (83.5,-9) {$2$};
\node[lab] at (110,-9) {$c_2=1$};
\draw[VR,line width=.6pt,decorate,decoration={brace,amplitude=2mm,mirror}] (31,-13) -- (81,-13);
\node[text=VR] at (56,-19.5) {精确还原$j$，不改变失真};
\end{tikzpicture}
